%% ma: L11-B27-0010
%% lop: 11
%% chuong: 7
%% bai: 27
%% dang: 11.27.7
%% loai: DS
%% muc: [TH, VD, VD, TH]
%% dap_an: "ĐSSĐ"
%% nguon: "(NBV)-ĐỀ SỐ 6-MỖI NGÀY 1 ĐỀ THI 2025.pdf (D06, MỖI NGÀY 1 ĐỀ - Đề số 6), Phần II Câu 1 (Chuyên Vĩnh Phúc 2025)"
%% kien_thuc: "Bài 25 (hai mặt phẳng vuông góc), Bài 26 (khoảng cách đến mặt phẳng), Bài 27 (thể tích khối chóp)"
%% trang_thai: cho-duyet
%% ly_do: "Dạng toán mới đề xuất 11.27.7 (hình chóp có mặt bên là tam giác đều vuông góc với đáy; Đúng/Sai về vuông góc, thể tích, khoảng cách), chờ giáo viên duyệt"
%% ghi_chu: "Đáp án gốc Đ S S Đ đúng (kiểm bằng tọa độ). Đề gốc không có hình"
\begin{ex}
Cho hình chóp $S.ABC$ có mặt bên $(SAB)$ vuông góc với mặt phẳng đáy và tam giác $SAB$ đều cạnh $2a$. Biết tam giác $ABC$ vuông tại $C$ và cạnh $AC=a\sqrt3$. Gọi $H$ là trung điểm $AB$.
\choiceTF
{\True Mặt phẳng $(SHC)$ và $(ABC)$ vuông góc với nhau.}
{Thể tích của khối chóp $S.ABC$ bằng $\dfrac{a^3}{6}$.}
{$d\big(C,(SAB)\big)=\dfrac{a\sqrt3}{3}$.}
{\True $SH\perp(ABC)$.}
\loigiai{Tam giác $SAB$ đều nên $SH\perp AB$ và $SH=a\sqrt3$. Vì $(SAB)\perp(ABC)$ theo giao tuyến $AB$ nên $SH\perp(ABC)$.

a) $(SHC)$ chứa $SH\perp(ABC)$ nên $(SHC)\perp(ABC)$. Đúng.

b) $BC=\sqrt{AB^2-AC^2}=a$, $S_{ABC}=\dfrac12\cdot a\sqrt3\cdot a=\dfrac{a^2\sqrt3}{2}$, $V=\dfrac13\cdot\dfrac{a^2\sqrt3}{2}\cdot a\sqrt3=\dfrac{a^3}{2}\ne\dfrac{a^3}{6}$. Sai.

c) $S_{SAB}=\dfrac{(2a)^2\sqrt3}{4}=a^2\sqrt3$, $d\big(C,(SAB)\big)=\dfrac{3V}{S_{SAB}}=\dfrac{3a^3/2}{a^2\sqrt3}=\dfrac{a\sqrt3}{2}\ne\dfrac{a\sqrt3}{3}$. Sai.

d) Đúng (chứng minh ở trên).}
\end{ex}
